\subsection{Proposition 2.7.2 of the master prompt: grounded behaviours and lower inputs}
\label{proof:P8}
\textbf{Setting and assumptions.} Fix a behaviour $b$, its concept graph
$\G$ and a set of principles $\Pi\subseteq\Pi_\tau$. Let $T$ be the training
inputs and $x^\ast$ a new input. Write $\mathrm{Desc}(\Pi)$ for the nodes of
$D_b(x^\ast)$ that are strictly downstream of $\Pi$ on a path to the output
node $o$ (including $o$).
\begin{enumerate}[label=(A\arabic*)]
\item \emph{Grounding on $x^\ast$.} $\Pi$ is a vertex cut of $D_b(x^\ast)$
between the input nodes and $o$.
\item \emph{Value locality downstream of the cut.} For every
$C\in\mathrm{Desc}(\Pi)$ and every $x\in T\cup\{x^\ast\}$,
$C(x)=\psi_C(\pa_{x}(C)(x))$, where $\pa_x(C)$ are the parents of $C$ that are
active on $x$, and $\psi_C$ is a fixed function (the system's derivation map,
which the exact derivation and interchange conditions identify on the
parent tuples seen in training).
\item \emph{$x^\ast$ is lower for $\Pi$.} There is a training input
$x_t\in T$ with $D_b(x_t)=D_b(x^\ast)$ on $\mathrm{Desc}(\Pi)\cup\Pi$ and
$\Pi(x_t)=\Pi(x^\ast)$ (the joint tuple of principle values is in range).
\end{enumerate}

\textbf{Proposition P8.} Under (A1)--(A3), $f_b(x^\ast)=f_b(x_t)$. Consequently, if
training error is zero on $x_t$ and the task's target factors through the
principle values ($\mathrm{target}(x)=\rho(\Pi(x))$ for some $\rho$ on
$T\cup\{x^\ast\}$), the system answers $x^\ast$ correctly.

\textbf{Proof.} Order $\mathrm{Desc}(\Pi)$ topologically. Claim: every
$C\in\mathrm{Desc}(\Pi)$ has $C(x^\ast)=C(x_t)$. Take $C$ and suppose the
claim holds for all earlier members. Each active parent $P$ of $C$ on $x^\ast$
lies on an active path to $o$ through $C$. If $P$ were neither in $\Pi$ nor in
$\mathrm{Desc}(\Pi)$, the active path from an input to $P$ (which exists
because $P\in D_b(x^\ast)$ lies on an input-to-output path) followed by
$P\to C\to\dots\to o$ would avoid $\Pi$: it cannot meet $\Pi$ after $C$
because $C$ is already downstream of $\Pi$ and the graph is acyclic, and it
does not meet $\Pi$ before $P$ by the choice of $P$. That contradicts (A1).
So every active parent of $C$ is in $\Pi$ or is an earlier member of
$\mathrm{Desc}(\Pi)$; by (A3) and the induction hypothesis its value on
$x^\ast$ equals its value on $x_t$, and by (A3) the active parent sets
agree. By (A2), $C(x^\ast)=\psi_C(\cdot)=C(x_t)$. Applying the claim to
$C=o$ gives $f_b(x^\ast)=f_b(x_t)$. If $f_b(x_t)$ is correct and the target
factors through $\Pi$, then $\mathrm{target}(x^\ast)=\rho(\Pi(x^\ast))=
\rho(\Pi(x_t))=\mathrm{target}(x_t)=f_b(x_t)=f_b(x^\ast)$. \hfill$\square$

\textbf{What the assumptions cost.} (A2) is the substantive one: it says the
nodes below the principles are computed from their parents' values only,
also on $x^\ast$. It is measurable on $x^\ast$ by within-fibre swaps and by
edge interventions from training partners (the interchange identification of
Proposition~\ref{prop:unseen}(b)). (A3) is the definition of a lower input,
restricted to the cut and below it; nodes upstream of the cut may be out of
range. The factorisation of the target is a property of the task and the
principles, not of the system. It fails exactly when a principle is a
shortcut (a high-generality node that predicts the target on the training
distribution only), and the proposition then predicts confident errors on
lower inputs, the K6 shortcut signature.

\textbf{Measurable prediction.} Among held-out inputs that are lower for the
principle cut, errors occur only where (A2) fails on the input (detected by
interventions on that input) or where the target does not factor through the
principles; errors on lower inputs with (A2) verified are therefore
concentrated in behaviours whose principles are known shortcuts. Tested in
K3 (moves) and K6 (shortcuts).
